\documentclass[11pt]{article}
\pdfmapfile{=cm.map}
\pdfmapfile{+cmextra.map}
\pdfmapfile{+symbols.map}
\pdfmapline{+tcrm1095 CMR10 <cmr10.pfb}
\usepackage[margin=1in]{geometry}
\usepackage{amsmath,amssymb,amsthm,mathtools}
\usepackage{booktabs}
\usepackage{hyperref}
\usepackage{listings}
\usepackage[numbers,sort&compress]{natbib}
\hypersetup{hidelinks}
\emergencystretch=3em

\newtheorem{theorem}{Theorem}[section]
\newtheorem{proposition}[theorem]{Proposition}
\newtheorem{lemma}[theorem]{Lemma}
\newtheorem{corollary}[theorem]{Corollary}
\theoremstyle{definition}
\newtheorem{definition}[theorem]{Definition}
\newtheorem{example}[theorem]{Example}
\theoremstyle{remark}
\newtheorem{remark}[theorem]{Remark}
\newtheorem{conjecture}[theorem]{Conjecture}

\newcommand{\Pow}{\mathcal P}
\newcommand{\Hyp}{\Omega}
\newcommand{\Feas}{\mathcal N}
\newcommand{\Docs}{\mathsf D}
\newcommand{\Claims}{\mathsf C}
\newcommand{\Frames}{\mathsf F}
\newcommand{\Constraints}{\mathsf K}
\newcommand{\Must}{\mathsf{Must}}
\newcommand{\May}{\mathsf{May}}
\newcommand{\Cannot}{\mathsf{Cannot}}
\newcommand{\Uncertain}{\mathsf U}

\lstset{basicstyle=\ttfamily\small,columns=fullflexible,breaklines=true,
  frame=single,framesep=4pt,xleftmargin=4pt}

\title{Dynamic Set-Theoretic Normalization of\\
Provenance-Bearing Hyper-Relational Frames}
\author{Project \path{thejeb}}
\date{July 2026}

\begin{document}
\maketitle

\begin{abstract}
We develop a deterministic semantics for normalizing claims and event frames
across a mutable document corpus.  The extraction front end may propose several
frame interpretations for each claim occurrence, and coreference is not forced
to one early clustering.  Instead, every exact interpretation and strict
coreference partition is a candidate normalization hypothesis.  Constraints
denote property sets of hypotheses, and the solver is literally set theoretic
estimation (STE): it repeatedly intersects its current state with property sets
and records the exact set reduction.  The resulting feasible family supports a
possible-worlds semantics for identity.  Coreference holding in every feasible
hypothesis defines a safe quotient; disagreement between feasible merge and
split hypotheses defines an explicit uncertainty overlay $\Uncertain$.

We prove document-insertion antitonicity, deletion expansion, insertion order
independence, a delta decomposition by newly activated provenance constraints,
identity and composition laws for feasible-world restriction, equivalence of
must-coreference, and exact status separation into inconsistent, must, cannot,
and uncertain cases.  A finite countermodel proves that possible coreference
and the uncertainty overlay need not be transitive even though every exact
hypothesis is a partition.  We prove existence of finite minimal document
supports for must- and cannot-coreference conclusions.  We then resolve six
initial conjectures under explicit, necessary qualifications: partial
variable-universe restriction; gluing by a cross-constraint obstruction with
at most one extension; a finite exact solver with tight $2^N$ state-space
bound; audited typed-relation lifting; exact marginal commutation for dependent
truth fibers; and occurrence-set identity with relational lineage.  The
accompanying Lean~4 development names every stated formal result.
\end{abstract}

\section{Introduction}

Cross-document event coreference is usually presented as clustering: a learned
model scores pairs or clusters, and an agglomerative or graph algorithm selects
one partition.  This is an effective predictive architecture, but it is an
awkward database semantics.  A database receives and retracts documents.  A
new source can distinguish events previously thought identical; deletion can
remove the only evidence that forced a merge or split.  Anaphora and omitted
arguments can remain genuinely unresolved.  A destructive union operation has
neither the provenance nor the inverse required by this setting.

Event identity itself is richer than lexical similarity.  Work on
quasi-identity and dense event annotation distinguishes strict identity from
subevent, near-identity, and related-event phenomena
\citep{hovy2013events,asr2014conceptual,pratapa2021dense}.  Recent resources
make decontextualization, difficult identity boundaries, and framing divergence
increasingly explicit \citep{zhao2025beyond,zhao2025framing}.  The correct
deterministic reaction is not to pretend every ambiguity has a unique answer.
It is to retain the admissible answers and state exactly what follows in all,
some, or none of them.

This paper formalizes that reaction.  We assume a high-recall extraction front
end.  Each immutable claim occurrence has source provenance and a nonempty set
of candidate hyper-relational frames.  An exact normalization hypothesis
chooses one candidate per claim and supplies an equivalence relation of strict
frame identity.  Typed nonidentity links---subevent, causation, temporal
relation, repeated instance, framing divergence---are conceptually separate
and do not participate in the identity quotient.

The contributions are:
\begin{enumerate}
  \item an STE solver semantics in which property-set application and exact
        reduction are the complete operational core;
  \item a provenance-indexed corpus semantics supporting insertion and
        deletion without destructive clustering;
  \item a must-coreference quotient and a nontransitive uncertainty overlay;
  \item a finite countermodel ruling out possible-coreference clustering;
  \item existence of minimal finite document certificates for settled links;
  \item contravariant feasible-world restriction and covariant canonical-frame
        maps, exposing the fundamental insertion/deletion asymmetry;
  \item a dependent lift of downstream truth estimation over all feasible
        normalizations.
\end{enumerate}

\section{The normalization universe}

Fix types of documents, claim occurrences, structured frames, exact
normalization hypotheses, and constraints:
\[
  \Docs,\qquad \Claims,\qquad \Frames,\qquad
  \Hyp,\qquad \Constraints.
\]
The Lean structure implementing these data is
\path{STE.DynamicFrame.Model}.

Each claim $c$ has an immutable source document $\operatorname{owner}(c)$ and a
nonempty candidate set $\Phi(c)\subseteq\Frames$.  Each exact hypothesis
$\omega\in\Hyp$ chooses
\[
  \operatorname{interpret}_{\omega}(c)\in\Phi(c)
\]
and supplies a strict-coreference relation $\sim_{\omega}$ on claims.  The
model requires $\sim_{\omega}$ to be an equivalence relation.  Thus every exact
hypothesis is an ordinary, logically valid partition.  Ambiguity is represented
between hypotheses, not by allowing one hypothesis to contain an invalid
overlapping ``partition.''

Every constraint $k\in\Constraints$ has two denotations.  First, its provenance
support is a set of documents
\[
  \operatorname{supp}(k)\subseteq\Docs.
\]
Second, it denotes an STE property set
\[
  S_k=\{\omega\in\Hyp\mid \omega\text{ satisfies }k\}.
\]
For an active corpus $D\subseteq\Docs$, the active constraints and feasible
normalizations are
\begin{align}
  K(D) &= \{k\mid \operatorname{supp}(k)\subseteq D\},\\
  \Feas(D) &= \bigcap_{k\in K(D)} S_k.
\end{align}
Empty-support logical or ontological constraints are active in every corpus.
A constraint derived jointly from several documents is active only while every
document in its support remains active.

This fixed ambient universe is deliberate.  Documents, claims, constraints,
and hypotheses are immutable mathematical objects.  Database mutation changes
an activation set.  A later variable-universe theory may garbage-collect or
forget inactive claims through explicit restriction maps, but semantic
retraction does not require erasing the objects that explain prior results.

\section{The solver is STE}

The phrase ``constraint solver'' can suggest an additional optimization or
clustering semantics.  There is none here.  The solver state is a set
$X\subseteq\Hyp$.  Applying $k$ is literal STE property-set application:
\begin{equation}
  \operatorname{apply}(X,k)=X\cap S_k.                 \label{eq:apply}
\end{equation}
The exact reduction caused by $k$ is
\begin{equation}
  \operatorname{red}(X,k)=X\setminus(X\cap S_k)
    =\{\omega\in X\mid \omega\notin S_k\}.           \label{eq:reduction}
\end{equation}
These are \path{Model.applyConstraint} and
\path{Model.solverReduction}.

\begin{proposition}[One-step STE laws]
Property application is contractive and idempotent; any two property
applications commute.  Moreover,
\[
  \operatorname{apply}(X,k)\cup\operatorname{red}(X,k)=X,
\]
and $\operatorname{red}(X,k)=\varnothing$ exactly when $X\subseteq S_k$.
\end{proposition}
\noindent Lean:
\path{applyConstraint_subset},
\path{applyConstraint_idempotent},
\path{applyConstraint_comm},
\path{apply_union_reduction}, and
\path{solverReduction_eq_empty_iff}.

For a finite schedule $[k_1,\ldots,k_n]$, define a run by repeated
application, starting from $X$, and define solving to start from $\Hyp$.

\begin{theorem}[Operational--denotational correspondence]
For every finite list $L$,
\begin{align}
  \omega\in\operatorname{run}(L,X)
    &\Longleftrightarrow
      \omega\in X\ \land\ \forall k\in L,\ \omega\in S_k,\\
  \operatorname{solve}(L)
    &=\bigcap_{k\in L}S_k.
\end{align}
Consequently, schedules with the same constraint membership have equal output;
in particular, permutations and duplicates do not matter.
\end{theorem}
\noindent Lean: \path{mem_run},
\path{run_eq_inter_partial},
\path{solve_eq_partialFeasibilitySet},
\path{solve_congr}, and \path{solve_perm}.

If a list enumerates exactly $K(D)$, the final operational state equals
$\Feas(D)$ (\path{solve_eq_feasible}).  Thus a SAT, SMT, ASP, BDD, or
custom partition engine is correct only insofar as it implements this meet and
its queries.  Such technologies are representations and acceleration
strategies, not alternate semantics.

Retraction follows the same point.  The system does not attempt to invert an
old merge.  It changes $K(D)$ using provenance and evaluates the new meet.  An
incremental implementation may reuse caches or truth-maintenance labels, but
its extensional correctness criterion is equality with the fresh STE meet.
This connects directly to truth-maintenance systems
\citep{doyle1979tms,dekleer1986atms}, provenance
\citep{green2007provenance}, and incremental view maintenance
\citep{budiu2023dbsp}.

\section{Dynamic corpus laws}

Let $D\subseteq E$.  Every constraint active in $D$ is active in $E$, hence
\begin{theorem}[Information antitonicity]
\[
  D\subseteq E\quad\Longrightarrow\quad \Feas(E)\subseteq\Feas(D).
\]
\end{theorem}
\noindent Lean: \path{Model.feasible_antitone}.  Read forward, document
addition removes hypotheses.  Read backward, removal can restore hypotheses.
The latter is \path{feasible_subset_after_remove}.

The elementary database laws are exact:
\begin{itemize}
  \item duplicate insertion is idempotent
        (\path{feasible_insert_idempotent});
  \item insertion order is irrelevant
        (\path{feasible_insert_comm});
  \item if $d\notin D$, adding and then removing $d$ restores precisely
        $\Feas(D)$ (\path{feasible_insert_then_remove}).
\end{itemize}

Define the constraint delta
\[
  \Delta K(D,E)=K(E)\setminus K(D).
\]
Then addition has an exact incremental decomposition.

\begin{theorem}[Delta decomposition]
If $D\subseteq E$, then
\[
  \omega\in\Feas(E)
  \Longleftrightarrow
  \omega\in\Feas(D)\ \land\
  \forall k\in\Delta K(D,E),\ \omega\in S_k.
\]
\end{theorem}
\noindent Lean: \path{mem_feasible_extension}.  If the delta is empty,
the feasible sets are equal
(\path{feasible_eq_of_no_new_constraints}).

Constraint activation interacts asymmetrically with union and intersection:
\[
 K(D\cap E)=K(D)\cap K(E),
 \qquad K(D)\cup K(E)\subseteq K(D\cup E).
\]
The second inclusion can be strict because $D\cup E$ may activate a constraint
whose support draws documents from both sides.  Lean:
\path{activeConstraints_inter} and
\path{activeConstraints_union_subset}.  This is the first genuinely
cross-document effect in the theory.

\section{Must, may, cannot, and uncertainty}

For $c,d\in\Claims$ and a consistent corpus, define
\begin{align}
 \Must_D(c,d)
   &\Longleftrightarrow \forall\omega\in\Feas(D),\ c\sim_\omega d,\\
 \May_D(c,d)
   &\Longleftrightarrow \exists\omega\in\Feas(D),\ c\sim_\omega d,\\
 \Cannot_D(c,d)
   &\Longleftrightarrow \forall\omega\in\Feas(D),\ c\not\sim_\omega d,\\
 \Uncertain_D(c,d)
   &\Longleftrightarrow
     (\exists\omega\in\Feas(D),\ c\sim_\omega d)\land
     (\exists\omega\in\Feas(D),\ c\not\sim_\omega d).
\end{align}
This is the certain/possible distinction of incomplete-database possible
worlds \citep{imielinski1984incomplete} applied to normalization partitions.

Consistency must be checked first.  If $\Feas(D)=\varnothing$, both universal
statements are vacuously true.  The implemented status type therefore has four
constructors:
\[
 \mathsf{inconsistent},\quad \mathsf{must},\quad
 \mathsf{cannot},\quad \mathsf{uncertain}.
\]
Lean: \path{Model.CorefStatus} and \path{Model.corefStatus}, with one
correctness theorem for each constructor.

\begin{theorem}[Safe quotient]
$\Must_D$ is an equivalence relation.  Therefore the deterministic normalized
frame collection
\[
  F'_D=\Claims(D)/\Must_D
\]
is well defined on active claims.
\end{theorem}
\begin{proof}
Reflexivity, symmetry, and transitivity hold in every exact relation
$\sim_\omega$ and are preserved by intersection over $\omega\in\Feas(D)$.
\end{proof}
\noindent Lean: \path{mustSame_equivalence},
\path{activeMustSetoid}, and \path{CanonicalFrame}.

Possible coreference is reflexive on a consistent corpus and symmetric, but no
general transitivity theorem is available.  $\Uncertain_D$ is symmetric and
irreflexive.  Must and cannot are each disjoint from uncertainty.  Lean:
\path{maySame_refl}, \path{maySame_symm},
\path{uncertain_symm}, \path{uncertain_irrefl},
\path{must_not_uncertain}, and
\path{cannot_not_uncertain}.

\subsection{A finite countermodel}

Let the mentions be $a,b,c$ and let two exact worlds be feasible.  The left
world has blocks $\{a,b\},\{c\}$; the right world has blocks
$\{a\},\{b,c\}$.  Each world is an ordinary partition.  Nevertheless,
\[
  \May(a,b),\qquad \May(b,c),\qquad \neg\May(a,c).
\]
Moreover $(a,b)$ and $(b,c)$ are uncertain while $(a,c)$ is cannot-corefer.

\begin{theorem}[Nontransitivity]
Neither $\May$ nor $\Uncertain$ is transitive in general.
\end{theorem}
\noindent The model and proofs are fully finite and constructive in
\path{Ste.DynamicFrame.Counterexample}; the theorem names are
\path{maySame_not_transitive} and
\path{uncertain_not_transitive}.

This establishes a design constraint.  $\Uncertain$ is an overlay between
must-classes, not a special kind of class.  Taking connected components of
possible or uncertain edges invents identifications not licensed in any exact
world.

\section{Restriction, quotient maps, and retraction asymmetry}

On the fixed ambient hypothesis universe, inclusion $D\subseteq E$ induces
\[
  \rho_{E,D}:\Feas(E)\longrightarrow\Feas(D),\qquad
  \rho_{E,D}(\omega)=\omega.
\]
This is a genuine restriction map: it forgets the additional requirements, not
the underlying hypothesis.

\begin{theorem}[Restriction laws]
Restriction by identity is identity, and for $D\subseteq E\subseteq F$,
\[
  \rho_{E,D}\circ\rho_{F,E}=\rho_{F,D}.
\]
A point of $\Feas(D)$ lies in the image of $\rho_{E,D}$ exactly when the same
ambient hypothesis is feasible in $E$.
\end{theorem}
\noindent Lean: \path{restrictFeasible_id},
\path{restrictFeasible_comp}, and
\path{mem_range_restrictFeasible_iff}.

Thus $D\mapsto\Feas(D)$ is already a contravariant set-valued construction on
the poset of active document sets.  Calling it a presheaf becomes literal once
the corpus-dependent universe and a target category are fixed.  Calling it a
sheaf would require a further gluing theorem; that is an open problem, not
terminology assumed here.

The quotient moves the other way.  Since
\[
  \Must_D(c,d)\Longrightarrow\Must_E(c,d),
\]
old must-classes map into new must-classes.  Active-claim inclusion and
relation preservation induce
\[
  q_{D,E}:F'_D\longrightarrow F'_E.
\]
Lean: \path{includeActiveClaim} and \path{canonicalMap}.
This map need not be injective: added evidence can eliminate all split worlds
and thereby merge old must-classes.  It also need not be surjective because
$E$ contains new active claims.

There is generally no quotient map in the deletion direction.  A pair forced
together in $E$ may become uncertain in $D$ after the supporting document is
removed.  Mapping one $E$-class back would then require splitting it, which is
not a function on quotient classes.  This is the formal reason that deletion
cannot be implemented as the inverse of an agglomerative merge history.

\section{Minimal document supports}

An explanation should identify documents sufficient for a settled conclusion.
For a finite document set $D$, say $D$ supports a must-link when $\Feas(D)$ is
nonempty and $\Must_D(c,d)$ holds.  Define cannot-support analogously.  A
support is minimal when no proper subset still supports the conclusion.

\begin{theorem}[Finite minimal certificates]
If a finite corpus supports a must-coreference conclusion, it contains an
inclusion-minimal finite support for that conclusion.  The same holds for a
cannot-coreference conclusion.
\end{theorem}
\begin{proof}
Among the finite supporting subsets of the original corpus, choose one of
minimum cardinality.  Every proper subset has smaller cardinality and therefore
cannot support the conclusion.
\end{proof}
\noindent Lean: the generic finite minimization theorem is
\path{exists_minimal_subfinset}; its two applications are
\path{exists_minimal_mustSupport} and
\path{exists_minimal_cannotSupport}.

The theorem is semantic, not an efficiency claim.  Computing all minimal
supports may be exponential.  ATMS labels, unsatisfiable cores, provenance
polynomials, and hitting-set dualization are candidate representations.  The
proof does establish a precise contract for an explanation engine: a returned
certificate should be sufficient, consistent, and inclusion-minimal.

\section{Lifting truth over normalization worlds}

Let $\mathsf T$ be a type of downstream truth worlds and let
\[
  T:\Hyp\longrightarrow\Pow(\mathsf T)
\]
assign a truth feasible set to each exact normalization.  The joint feasible
set is
\[
  J(D)=\{(\omega,t)\mid \omega\in\Feas(D),\ t\in T(\omega)\}.
\]
This is \path{TruthLift.jointFeasible}.  A property $P(\omega,t)$ is certain
when it holds throughout $J(D)$ and possible when it holds somewhere in
$J(D)$.

\begin{theorem}[Truth-lift monotonicity]
If $D\subseteq E$, every property certain over $J(D)$ is certain over $J(E)$,
and every property possible over $J(E)$ was possible over $J(D)$.
\end{theorem}
\noindent Lean: \path{certain_mono} and
\path{possible_antitone}.

Not every feasible normalization need support a downstream truth world.  The
projection of $J(D)$ onto normalization hypotheses is contained in
$\Feas(D)$, with equality if every feasible normalization has a nonempty truth
fiber.  Lean: \path{normalizationProjection_subset} and
\path{normalizationProjection_eq}.

The first commutation result covers the clean independence case.

\begin{theorem}[Normalization-independent commutation]
Assume $\Feas(D)\ne\varnothing$, $T(\omega)=T_0$ for every $\omega$, and
$P(\omega,t)=P_0(t)$.  Then
\[
  \forall(\omega,t)\in J(D),\ P_0(t)
  \Longleftrightarrow
  \forall t\in T_0,\ P_0(t).
\]
Existentially, joint possibility is equivalent to consistency of
$\Feas(D)$ together with ordinary possibility in $T_0$.
\end{theorem}
\noindent Lean: \path{certain_constant_iff} and
\path{possible_constant_iff}.

This identifies the exact modularity boundary.  Normalization-first inference
commutes automatically when truth semantics ignores normalization.  The
research problem is to weaken that assumption: characterize dependent
$T(\omega)$ and $P(\omega,t)$ for which certain answers can be computed without
enumerating every joint world.

\section{Status evolution under corpus change}

The modal laws give a concise transition discipline.  Under $D\subseteq E$:
\begin{align*}
 \Must_D(c,d)&\Longrightarrow\Must_E(c,d),\\
 \Cannot_D(c,d)&\Longrightarrow\Cannot_E(c,d),\\
 \May_E(c,d)&\Longrightarrow\May_D(c,d),\\
 \Uncertain_E(c,d)&\Longrightarrow\Uncertain_D(c,d).
\end{align*}
Thus addition can resolve uncertainty into must or cannot, or make the entire
problem inconsistent.  It cannot turn a consistent settled fact into
uncertainty within the hard-constraint semantics.  Deletion reverses the
information order: a settled link may become uncertain, but a possibility
cannot be invented by adding a constraint.

These are \path{mustSame_mono}, \path{cannotSame_mono},
\path{maySame_antitone}, and \path{uncertain_antitone}.  They are useful
both as theorems and as property-based tests for an incremental implementation.

\section{Relationship to learned coreference}

Nothing in the theory forbids neural models.  They can propose candidate frame
interpretations, candidate links, and likely high-reduction criteria.  They can
rank feasible worlds.  What they cannot do silently is turn an uncalibrated
score threshold into a semantic impossibility.

This yields three distinct interfaces:
\begin{enumerate}
  \item \emph{generation}: place defensible worlds inside the ambient candidate
        universe;
  \item \emph{hard admissibility}: apply audited property sets by STE;
  \item \emph{ranking}: order the worlds that remain.
\end{enumerate}
Candidate recall is a fairness condition on generation.  Hard-constraint
soundness is a fairness condition on property sets.  Ranking does not alter
must/cannot guarantees unless the ranked tail is explicitly discarded, in
which case that discard is itself a property-set application that must be
named and audited.

\section{Mechanization map}

The development is split by responsibility:
\begin{center}
\begin{tabular}{ll}
\toprule
Lean module & Formal content \\
\midrule
\path{Ste.DynamicFrame} & model, feasibility, modal relations, quotient \\
\path{Ste.DynamicFrame.Laws} & dynamic laws, restriction, status, quotient map \\
\path{Ste.DynamicFrame.Solver} & property application, reduction, finite runs \\
\path{Ste.DynamicFrame.Counterexample} & three-mention nontransitivity model \\
\path{Ste.DynamicFrame.Support} & finite minimal document certificates \\
\path{Ste.DynamicFrame.PresheafGluing} & partial restriction and exact gluing \\
\path{Ste.DynamicFrame.FiniteSolver} & executable solver and tight bounds \\
\path{Ste.DynamicFrame.RelationAlgebra} & audited typed necessary relations \\
\path{Ste.DynamicFrame.TruthLift} & dependent downstream feasible family \\
\path{Ste.DynamicFrame.Lineage} & occurrence-set identifiers and lineage \\
\bottomrule
\end{tabular}
\end{center}

The root module \path{Ste.lean} imports the complete theory.  The project is
pinned to Lean~4.32.0 and Mathlib~4.32.0.  The repository workflow builds the
root library on pushes to the development branch.  The source and this paper
use a strict convention: a named Lean theorem is the authoritative statement.
The next section replaces the earlier conjectures by qualified theorems and
records the boundary at which each claim becomes false.

\section{Resolution of the six conjectures}

The original conjectures mixed extensional existence claims with computational
ones.  None should be accepted without its domain assumptions.  This section
states what is actually achievable, proves it, and records sharp constants or
small counterexamples where they matter.

\subsection{Variable-universe restriction is necessarily partial}

Write $W(D)=\Feas(D)$ for the subtype of exact worlds feasible at corpus $D$.
For $D\subseteq E$, define
\[
  \rho_{E,D}:W(E)\longrightarrow W(D),\qquad \rho_{E,D}(h)=h.
\]
The earlier restriction proofs immediately give
\[
  \rho_{D,D}=\mathrm{id},\qquad
  \rho_{E,D}\circ\rho_{F,E}=\rho_{F,D}.
\]
Lean: \path{restrictWorld_id} and \path{restrictWorld_comp}.  Thus feasible
worlds form a contravariant functor on the document-set poset without any new
axiom.

The active claim universe behaves differently.  Insertion gives a total map
$\Claims(D)\to\Claims(E)$, but deletion cannot give a total reverse map: a claim
whose source lies in $E\setminus D$ has no active target.  Define a partial
restriction relation by equality of immutable occurrence identifiers.

\begin{theorem}[Partial claim restriction]
For $D\subseteq E$ and $c\in\Claims(E)$, a restriction of $c$ to
$\Claims(D)$ exists exactly when the source document of $c$ lies in $D$; when
it exists, it is unique.
\end{theorem}
\begin{proof}
If a restricted active claim is equal to $c$, its source is active in $D$.
Conversely, source membership constructs the subtype element directly.
Equality of two such subtype elements follows from equality of their immutable
occurrences.
\end{proof}
\noindent Lean: \path{claimRestriction_exists_iff} and
\path{claimRestriction_unique}.  The exact fiber constant is therefore
$0$ or $1$, not always $1$.  Local strict-coreference is preserved
definitionally by restricting the world and including the surviving claims
(\path{localSame_restrict}).  This is the precise variable-universe result;
a total deletion map was the false part of the unqualified conjecture.

\subsection{Gluing is controlled by cross-document constraints}

For an ambient hypothesis $h$, define the obstruction
\begin{align*}
 \mathcal O_{D,E}(h)=\{k\in\Constraints:
   &\operatorname{supp}(k)\subseteq D\cup E,\\
   &\operatorname{supp}(k)\nsubseteq D,
    \operatorname{supp}(k)\nsubseteq E,
    h\notin S_k\}.
\end{align*}
Only genuinely cross-document constraints occur in this set.

\begin{theorem}[Exact union obstruction]
For every $h$,
\[
 h\in\Feas(D\cup E)
 \Longleftrightarrow
 h\in\Feas(D)\land h\in\Feas(E)\land
 \mathcal O_{D,E}(h)=\varnothing.
\]
\end{theorem}
\begin{proof}
The forward direction follows by feasibility antitonicity and satisfaction of
every union-active constraint.  Conversely, a union-active constraint is
active on $D$, active on $E$, or active on neither side.  The first two cases
follow from local feasibility.  In the third case, violation would place the
constraint in $\mathcal O_{D,E}(h)$, contradicting emptiness.
\end{proof}
\noindent Lean: \path{mem_feasible_union_iff}.

Two local views are called compatible here when they expose the same ambient
exact hypothesis.  Under that qualification, the global-extension type is
nonempty exactly when the obstruction vanishes
(\path{globalExtension_nonempty_iff}), and it is a subsingleton
(\path{globalExtension_subsingleton}).  Hence the number of extensions is at
most $1$.  For finite $\Constraints$ with decidable support and satisfaction,
the obstruction is computable by one filter and contains at most
$|\Constraints|$ elements.  If compatibility is weakened to agreement only on
overlap, uniqueness is no longer promised: residual identity choices outside
the overlap are precisely additional global worlds.

\subsection{The finite exact solver and its tight limit}

Assume the ambient exact hypothesis universe is finite of cardinality
\[
  N=|\Hyp|,
\]
and satisfaction is decidable.  Represent a solver state by a finite set and
implement property application by filtering.

\begin{theorem}[Executable exactness]
Finite filtering denotes exactly abstract STE intersection, and a finite run
denotes exactly the corresponding abstract run and partial feasibility set.
\end{theorem}
\begin{proof}
Membership after filtering is membership in the input conjoined with
satisfaction.  Induction on the constraint list then identifies every finite
run with repeated property-set intersection.
\end{proof}
\noindent Lean: \path{coe_applyFinite}, \path{coe_runFinite}, and
\path{coe_solveFinite}.

Every nontrivial application removes at least one world
(\path{applyFinite_card_lt}).  Total reduction from any state is at most $N$
(\path{reductionCost_le_maxWorlds}).  These are useful operational constants:
the explicit representation stores at most $N$ survivors, a full must/may scan
examines at most $N$ worlds, and no execution can eliminate more than $N$
distinct worlds.

The unconditional compactness conjecture, however, is false.

\begin{theorem}[Tight extensional state bound]
There are exactly $2^N$ extensional solver states over $N$ ambient worlds, and
unrestricted crisp STE can realize every one of them with a single property
set.
\end{theorem}
\begin{proof}
The solver states are the powerset of the $N$-element universe, whose
cardinality is $2^N$.  Given any $A\subseteq\Hyp$, use one constraint with
property set $S_\star=A$; its feasibility set is exactly $A$.
\end{proof}
\noindent Lean: \path{card_allExactStates} and
\path{arbitrary_subset_as_single_property}.  Therefore no uniform
subexponential extensional representation follows from STE alone.  BDDs,
decomposable circuits, treewidth bounds, or factorized partition structures
remain worthwhile, but their compactness theorems must assume structure in the
generated property sets.

\subsection{Typed relations lift exactly when their rules are audited}

The mechanization gives each relation kind a source and target sort and keeps
strict identity, subevent, supersession, before, causes, repeated instance, and
framing divergence distinct.  A concrete world semantics supplies the laws it
actually validates: transitivity for subevent, supersession, and before;
causation implying temporal precedence; equivalence laws for repeated
instance; and symmetry plus irreflexivity for framing divergence.

\begin{theorem}[Necessary-relation lifting]
Let a composition rule
\[
 R_i(x,y)\land R_j(y,z)\Longrightarrow R_k(x,z)
\]
hold in every exact world.  Replacing each relation by its must-version---the
intersection over feasible worlds---preserves the same rule.
\end{theorem}
\begin{proof}
Fix a feasible world $h$.  Both must-premises hold in $h$, so world-level
soundness gives the conclusion in $h$.  Since $h$ was arbitrary, the
must-conclusion follows.
\end{proof}
\noindent Lean: \path{CompositionRule} and \path{must_compose}.  Specialized
proofs are \path{mustSubevent_trans}, \path{mustSupersession_trans},
\path{mustBefore_trans}, and \path{mustCausesBefore}; the repeated-instance and
framing-divergence laws are lifted separately.  The proof preserves one rule
application exactly: two necessary premises produce one necessary conclusion.
No theorem asserts, for example, transitivity of causation or framing
divergence.  Names do not license composition constants; audited world laws do.

Consistency is required to recover endpoint data from a necessary statement,
because universal relations over an empty feasible set are vacuous.  This is
explicit in \path{mustRel_typed}, \path{mustStrictIdentity_iff}, and
\path{not_mustDivergence_self}.

\subsection{Dependent truth commutes through its marginal union}

For dependent truth fibers $T(h)$ define
\[
  T_\exists(D)=\{t:\exists h\in\Feas(D),\ t\in T(h)\}.
\]
The fibers may vary arbitrarily.

\begin{theorem}[Marginal commutation]
For every truth-only predicate $Q$,
\begin{align*}
 \forall h\in\Feas(D)\ \forall t\in T(h),\ Q(t)
 &\Longleftrightarrow \forall t\in T_\exists(D),\ Q(t),\\
 \exists h\in\Feas(D)\ \exists t\in T(h),\ Q(t)
 &\Longleftrightarrow \exists t\in T_\exists(D),\ Q(t).
\end{align*}
\end{theorem}
\begin{proof}
Both equivalences follow by unpacking the existential witness in the definition
of $T_\exists(D)$.  No fiber constancy or nonemptiness premise is used.
\end{proof}
\noindent Lean: \path{certain_marginal_iff} and
\path{possible_marginal_iff}.

A genuinely dependent query $P(h,t)$ also commutes if, on jointly feasible
support, there is a $Q(t)$ with $P(h,t)\leftrightarrow Q(t)$.  This checkable
condition is \path{SupportedInvariant}; the exact universal and existential
results are \path{certain_supportedInvariant_iff} and
\path{possible_supportedInvariant_iff}.  The loss constant is zero: these are
logical equivalences, not approximations.  The condition is also meaningful.
Two feasible normalizations sharing one truth world but disagreeing on $P$
rule out every truth-only factor $Q$; this two-normalization, one-truth-world
witness is formalized by \path{no_supportedInvariant_of_disagreement}.

\subsection{Stable identity is a member set; history is a relation}

For a canonical must-class $q$ at snapshot $D$, define
\[
  \operatorname{id}_D(q)=
  \{c\in\Claims(D):c\text{ belongs to }q\}.
\]
These are immutable claim occurrence identifiers, not generated cluster names.

\begin{theorem}[Complete snapshot identifier]
Every $\operatorname{id}_D(q)$ is nonempty, and
$\operatorname{id}_D(q)=\operatorname{id}_D(r)$ implies $q=r$.
For $D\subseteq E$,
\[
 \operatorname{id}_D(q)\subseteq
 \operatorname{id}_E(q_{D,E}(q)).
\]
\end{theorem}
\begin{proof}
A representative belongs to its own equivalence class.  Equal member sets put
each representative in the other's class, so quotient soundness gives equal
classes.  Under insertion, claims remain active and must-coreference is
monotone, proving the inclusion.
\end{proof}
\noindent Lean: \path{canonicalMembers_nonempty},
\path{canonicalMembers_injective}, and
\path{canonicalMembers_subset_map}.  Thus the minimum identifier size is $1$,
and the identifier has zero collisions inside a snapshot.

Across arbitrary snapshots define lineage by nonempty intersection of member
sets.  It is reflexive for live frames, symmetric, and every frame has lineage
to its insertion image (\path{lineage_refl}, \path{lineage_symm}, and
\path{lineage_canonicalMap}).  It is intentionally not transitive: the sets
$\{a\}$, $\{a,b\}$, and $\{b\}$ give the smallest counterexample.  Lean proves
this on a two-element type in \path{overlap_not_transitive}.  Consequently a
single immutable scalar identifier cannot simultaneously model a merge and
both later split children without collision or arbitrary choice.  Stable
snapshot identity plus relational lineage is achievable; globally functional
identity through splits is not.

\subsection{Achievability constants}

\begin{center}
\begin{tabular}{p{0.27\textwidth}p{0.24\textwidth}p{0.37\textwidth}}
\toprule
Result & Constant & Qualification or sharp limit \\
\midrule
Claim restriction & fiber size $0$ or $1$ & source must remain active \\
Compatible gluing & extensions $\leq 1$ & locals name the same ambient world \\
Finite exact solver & $N$ survivors; $2^N$ states & exponential bound is tight \\
Typed composition & one sound rule step & only audited rules lift \\
Truth commutation & zero logical loss & invariant on feasible support \\
Snapshot identity & size $\geq1$; zero collisions & history is overlap, not a function \\
\bottomrule
\end{tabular}
\end{center}

Other engineering questions retain exact correctness oracles: incremental
output must equal a fresh STE meet; insertion order must not matter; deletion
must reactivate every world excluded only by removed support; and every
published settled link should admit a finite minimal certificate when the
active corpus is finite.

\section{Conclusion}

Dynamic frame normalization is best understood as set-theoretic estimation
over exact partitions, not as an irreversible clustering task.  Property-set
application is the solver.  Provenance selects the active meet.  Must
coreference is the intersection of feasible equivalence relations and supplies
the only generally safe deterministic quotient.  `May' and $\Uncertain$ are
modal summaries, and a finite machine-checked countermodel shows why they
cannot be collapsed into partitions.  Document addition and deletion act in
opposite information directions; feasible worlds restrict contravariantly,
while canonical quotients map covariantly under addition.  Minimal finite
supports provide an explanation contract, and the dependent truth lift keeps
normalization uncertainty intact downstream.

The result is a substantive mathematical kernel for a live document database:
ambiguity is represented, reduction is exact, retraction has a denotation,
cross-document gluing has an obstruction, and finite implementation has both
an executable correctness theorem and a tight worst-case bound.  The earlier
open conjectures are no longer slogans: each is either proved under explicit
hypotheses or delimited by a machine-checked counterexample.  The principal
remaining algorithmic question is structural rather than semantic---which
properties of real candidate graphs make their exact $2^N$ worst case compress
in practice while retaining deletion, explanation, and must/may queries.

\bibliographystyle{plainnat}
\bibliography{../../refs}

\end{document}
